% Канонический TikZ-рисунок варианта 11. Править только здесь.
\newcommand{\figIncenterBisectors}[1]{%
\begin{tikzpicture}[scale=0.75,line cap=round,line join=round,every node/.style={font=\small}]
  \coordinate (A) at (0,0); \coordinate (B) at (8,0); \coordinate (C) at (3.5,9);
  \coordinate (D) at (5.96,4.08); \coordinate (E) at (1.55,3.99); \coordinate (O) at (3.80,2.60);
  \draw[line width=0.9pt] (A)--(B)--(C)--cycle;
  \draw[line width=0.85pt] (A)--(D); \draw[line width=0.85pt] (B)--(E);
  \fill (O) circle (1.4pt);
  \pic[draw,line width=0.6pt,angle radius=6mm] {angle=B--A--D};
  \pic[draw,line width=0.6pt,angle radius=6mm] {angle=D--A--C};
  \pic[draw,line width=0.6pt,angle radius=5mm] {angle=C--B--E};
  \pic[draw,line width=0.6pt,angle radius=6.4mm] {angle=C--B--E};
  \pic[draw,line width=0.6pt,angle radius=5mm] {angle=E--B--A};
  \pic[draw,line width=0.6pt,angle radius=6.4mm] {angle=E--B--A};
  \pic[draw,line width=0.65pt,angle radius=7mm] {angle=A--C--B};
  \node at (3.55,7.15) {$#1^\circ$};
  \node[below left] at (A) {$A$}; \node[below right] at (B) {$B$}; \node[above] at (C) {$C$};
  \node[right] at (D) {$D$}; \node[left] at (E) {$E$}; \node[below] at (O) {$O$};
\end{tikzpicture}}
